\section{The conditional estimator at large counts}\label{sec:high}

The second moment just proved suffices when the mass of a cell
is small. At larger counts we will obtain a conditional variance
bound \(C n^\eta\log^3(n)/N\) on the good pilot event.
Fix \(\gamma>0\) and suppose
\(n^\gamma\le N\le n\) with \(w=w_I\) a probability density.
We keep \(t=\log n/N\), \(q=q_I\), \(e=\|w-q\|_2\) and
\(r_*=t/K^2\). Split the observations into three independent
fractions of sizes comparable to \(N\), for a coarse pilot,
the interval pilot, and the correction group.

The empirical ordinary projection from the coarse pilot,
with coefficients rounded at scale \(n^{-3}\), gives a
polynomial \(a\) of uniformly bounded coefficients. It satisfies
\begin{equation}\label{eq:coarse}
 \|a-q\|_\infty\le \delta=L_3\sqrt{\log n/N}
\end{equation}
with probability at least \(1-n^{-M}\) for any prescribed
\(M\) with \(L_3\) sufficiently large, by concentration
of the fixed bounded basis scores. For a fixed \(D_0>2\), set
\[
 S=\{x\in W: |a(x)|\le D_0\delta\},\qquad T=W\setminus S .
\]
Conditional on the coarse pilot, take intervals for the
ordinary scores and the finite dictionary of kernels obtained
by starting the cuts of Section~\ref{sec:angular} at \(S\).
This includes the indicator of \(S\).
The real interval pilot is independent of \(a\) and the
correction sample. On \eqref{eq:coarse},
\(\sup_S|q|\le(D_0+1)\delta\) and \(|S|\le C\delta^\beta\),
by Lemma~\ref{lem:sublevel}. If \(\inf_W|q|\ge1/2\),
then \(S\) is empty for \(n\) large enough, uniformly in
\(n^\gamma\le N\le n\).

On \(S\) include the linear term \(\int_S q\log\delta\).
Include the nonlinear correction from the angular construction
only if \(|S|>0\) and the interval pilot's mean
\(\widehat p_S\) of \(\one_S\) satisfies
\(\widehat p_S\le L_4\delta\). On valid pilots, retention implies
\(p_S\le C\delta\), and rejection of a positive-volume \(S\)
implies \(p_S\ge (L_4/2)\delta\) for sufficiently large \(n\).
Both implications follow from
\(|\widehat p_S-p_S|\le C(\sqrt{tp_S}+t)\) with \(t=o(\delta)\).
Choose \(L_4/2>D_0+2\).
In the rejected case,
\[
 \sqrt{|S|}\,e\ \ge\ \int_S(w-q)
   \ \ge\ p_S-(D_0+1)\delta |S|\ \ge\ c\delta .
\]
Omitting the centered nonlinear term then incurs bias at
most \(C\delta|S|\le Ce\), since \(|q|\le(D_0+1)\delta\)
on \(S\). A null set contributes zero.

For a retained \(S\), start with both support and evaluation
region equal to \(S\), with its weighted kernel and initial
mass \(p_S\). In \eqref{eq:piece} the ordinary projection
\(q\) in the linear term still uses all of \(\Omega\).
Use \(a_*=\delta\), \(r_*=t/K^2\) and the same Bernstein
intervals throughout the \(J\) selections. The stage bias and mass
recurrence from Section~\ref{sec:angular} give
\(C r_*(1+(p_S/r_*)^\epsilon)+Ce\) for the filter.
The replacement of regional projections by \(q\) is needed
from the first weighted stage. At each stage the estimate is
\(C r_*+C\log n\sqrt{|S|}\,e\), by \eqref{eq:regional},
\(\|q\|_\infty\le C\), and \(r_*\gtrsim (n\log n)^{-1}\).
It is \(O(r_*+e)\) because \(|S|\le C\delta^\beta\),
whose positive power dominates the logarithm for
\(N\ge n^\gamma\). Thus the centered angular part, with the
gate included, has absolute conditional bias at most
\begin{equation}\label{eq:high-bias}
 C\bigl[e+r_*(1+r_*^{-\epsilon}
                        \one_{\{\inf_W|q|<1/2\}})\bigr],
\end{equation}
using \(p_S\le1\) and the empty-set assertion.

On valid retained histories, all descendant restricted-score
intervals have endpoints \(O(\delta)\), since their means
are bounded by \(p_S\le C\delta\) and their widths are
\(O(\sqrt{t p_S}+t)\). Thus the signed filter acts on
\(G(c)\) for \(|c|\le C\delta\) on the rectangles, giving
coefficient norm at most \(C^K\delta\log n\).
The global score intervals have widths \(O(\delta)\).
The ordinary projection evaluated at their lower endpoints
differs from \(q\) by \(O(\delta)\) uniformly, by the bounded
basis and fixed dimension. Its value on \(S\) is \(O(\delta)\).
Thus the constant term and each coefficient of that
projection as an affine expression in normalized scores is
\(O(\delta)\) on \(S\). Integrating its product with
\(\log(\ell_j r_*/\delta)\) on any descendant preserves a
coefficient norm bound \(C\delta\log n\), by the log-envelope
integral in Section~\ref{sec:angular}. Sum over the finite
paths via the outgoing probability norm bound. The factorial
moment argument now bounds the conditional second moment of
the centered angular estimate by \(C_J^K\delta^2\log^2 n\).
Its absolute conditional mean is at most \(C\delta\log n\):
at the true rectangle coordinates, the absolute filtered
value is bounded by \(C\delta\log n\) by the uniform
\(L^1\) norm of the filter, and the path weights form a
probability distribution with fixed depth.

On \(T\), put \(z=(q-a)/a\). On \eqref{eq:coarse} we have
\(|z|\le 1/D_0\). The pointwise identity
\[
 \psi(q)=q\log|a|+a(1+z)\log(1+z)
\]
allows the use of the Taylor polynomial of
\((1+z)\log(1+z)\) at zero of degree
\(M_T=\lceil C_T\log n\rceil\). Its nonconstant coefficients
have magnitude at most one. The integrated error is bounded by
\(C D_0^{-M_T}\int_T |a|\), which is \(o(r_*)\) for fixed
large \(C_T,D_0\). Estimate each power of \(z(x)\) using
ordered products of the score
\[
 Y(x,X)=\frac{K_\Omega(x,X)-a(x)}{a(x)}
\]
on the correction sample, with its actual falling factorial
denominator. Since \(\E_w K_\Omega(x,X)=q(x)\), the mean
of this signed score is \(z(x)\). On \(T\), the denominator
\(a(x)\) is separated
from zero by the observed threshold on every pilot history.
On \eqref{eq:coarse}, \(|\E Y|\le1/D_0\) and
\(\E Y^2\le C(1+N/(D_0^2\log n))\).
The overlap count from Section~\ref{sec:angular} now retains the
unshared means. For an order-\(k\) statistic, also with one
observation omitted, it bounds the second moment by
\[
 \bigl(D_0^{-2}+Ck(1+N/(D_0^2\log n))/N\bigr)^k .
\]
For \(k\le M_T\) this is at most \(\rho^{2k}\) with
a fixed \(\rho<1\), by choosing \(D_0\) sufficiently large
after \(C_T\), and then \(n\) large. The correction size
is of order \(N\) and \(M_T=o(N)\).

To bound the variance of the integrated Taylor statistic,
replace one correction observation \(X_i\) by an independent
copy \(X_i'\). At order \(k\) the difference of the ordered products statistic equals
\((k/m_c)(Y(x,X_i)-Y(x,X_i'))\) times the order-\((k-1)\)
statistic omitting \(i\). Here \(k/m_c=k(m_c-1)_{k-1}/(m_c)_k\).
The factors are independent conditional on the pilots, and
\[
 a(x)\bigl(Y(x,X_i)-Y(x,X_i')\bigr)
  = K_\Omega(x,X_i)-K_\Omega(x,X_i') .
\]
Thus after multiplication by \(a(x)\) the first factor has
\(L^2\) norm at most \(C k/m_c\)
by the fixed kernel bound. The other factor has \(L^2\) norm
at most \(\rho^{k-1}\). Minkowski for the integral and
Taylor sum gives \(O(1/m_c)\) for the \(L^2\) effect of
each replacement. The Efron--Stein variance bound (the sum
of these squared replacement effects divided by two) gives
variance \(O(1/N)\).

Estimate the linear contribution \(\int_T q\log|a|+
\int_S q\log\delta\) by averaging the corresponding
kernel integral on the correction sample. The score is bounded
by \(C\log n\) from the threshold on \(T\) and the uniform
bound on \(a\). This term has conditional variance
\(O(\log^2 n/N)\). The three parts may share the correction
observations, so add their conditional standard deviations.
The resulting statistic estimates \(\int_W\psi(q)\) with
bias bounded by \eqref{eq:high-bias} and variance
\begin{equation}\label{eq:high-var}
 C n^\eta\log^3 n/N
\end{equation}
on the stated good actual pilots, after decreasing \(c_K\)
as necessary. Its conditional mean is bounded by \(C\log n\):
use the retained angular mean bound above, the uniform Taylor
mean bound from the geometric series at \(|z|\le 1/D_0\),
and the bounded linear score. The estimator is defined by
observed conditions, and the probability of failure of the
analytic pilot bounds is \(O(n^{-M})\) conditional on \(N\).
